\documentclass[11pt, letterpaper]{article}

\usepackage[
    top=0.7in,
    bottom=0.7in,
    left=0.85in,
    right=0.85in,
]{geometry}
\usepackage{charter}
\usepackage{microtype}
\usepackage[T1]{fontenc}
\usepackage[utf8]{inputenc}
\usepackage{hyperref}
\hypersetup{
    colorlinks=true,
    urlcolor=black,
    linkcolor=black,
}

\setlength{\parskip}{0.45em}
\setlength{\parindent}{0pt}
\linespread{1.0}

\pagestyle{empty}

\begin{document}

\begin{flushright}
Haoming Wang \,$|$\, 334 N Craig, Pittsburgh, PA \,$|$\, \href{mailto:haw200@pitt.edu}{haw200@pitt.edu} \,$|$\, \today
\end{flushright}

\vspace{-0.3em}

\noindent\textbf{To the MLSys 2026 Travel Grant Selection Committee}

\vspace{0.2em}

\noindent Dear Committee Members,

I am writing to apply for the MLSys 2026 Student Travel Grant. I am a fourth-year Ph.D.\ candidate in Electrical and Computer Engineering at the University of Pittsburgh, advised by Prof.\ Wei Gao in the Intelligent Systems Lab. My research sits squarely at the intersection that MLSys exists to advance --- building the systems that make modern ML models actually run, learn, and adapt on real hardware. I am U.S.-based (Pittsburgh, PA) and currently a degree-seeking student with an anticipated graduation in May 2027.

Over the past three years I have developed a sustained research agenda on \textbf{generative AI for mobile systems and on-device learning}, with first-authored papers at MobiCom 2025 (FedDC, 17.1\% accept), MobiSys 2025 (XPerT, 18.0\% accept), AAAI 2025, and CVPR 2026 (InfiniBench, Oral). FedDC introduced the first federated-learning technique to address the \emph{correlated} coupling of device delays and data heterogeneity in real AIoT deployments, using server-side gradient inversion to compensate stale updates and raising accuracy by up to 34\% across 13 heterogeneous IoT devices. XPerT expedites on-device LLM personalization by replacing from-scratch fine-tuning with explainable selection of cached personalized LLMs, cutting on-device fine-tuning time by 83\% and communication cost by 96.5\% on Pixel and OnePlus smartphones. My ongoing work --- MosaicThinker, a training-free pipeline that enables small VLMs to perform cross-frame spatial reasoning on Jetson Orin, AR glasses, and mobile phones --- extends this trajectory toward embodied AI on edge devices.

These threads are exactly the conversations MLSys convenes: efficient training and inference, federated and distributed learning, on-device adaptation, and the systems--algorithm co-design needed to make foundation models tractable beyond the data center. Attending MLSys 2026 in person would let me (1) present and discuss this line of work with the systems community whose feedback most directly shapes its trajectory, (2) absorb advances in compiler-, runtime-, and hardware-level techniques that I cannot fully grasp from papers alone, and (3) build collaborations with peers and senior researchers as I move into the final phase of my Ph.D.\ and prepare for the academic job market.

Travel funding is the practical obstacle. As a Ph.D.\ student supported primarily by a research assistantship, I do not have a discretionary travel budget, and my advisor's grant funds for FY2026 are committed to existing project obligations and student conference travel for our group's other accepted venues. A travel grant would meaningfully lower the barrier and allow me to participate fully in the conference rather than only its virtual components. I would be deeply grateful for the committee's consideration.

\vspace{0.4em}

\noindent Sincerely,

\vspace{0.6em}

\noindent Haoming Wang \\
Ph.D.\ Candidate, Electrical and Computer Engineering, University of Pittsburgh \\
\href{mailto:haw200@pitt.edu}{haw200@pitt.edu} \,$|$\, \href{https://haomingwang645.github.io/}{haomingwang645.github.io}

\end{document}
